
\paragraph{Statement A.4.2 (Proved reduction; its needed uniform estimate remains undecided).}
For each integer $\ell$ define a unitary character on the nonzero ideals
of $\mathbb Z[\sqrt3]$ by
\[
 \psi_\ell((\alpha))=
 \left|\frac{\alpha}{\alpha'}\right|^{i\pi\ell/L_0},
 \qquad a_\ell(m)=\sum_{N\mathfrak a=m}\psi_\ell(\mathfrak a),
\]
where the prime denotes the nontrivial real embedding. Then
\begin{equation}\label{A:twisted-shift}
 T_\ell(X):=\sum_{z\leq X}\sum_{N\mathfrak a=m_z}
 e^{2\pi i\ell\theta(\mathfrak a)}
 =(-1)^\ell\sum_{z\leq X}a_\ell(3z-1)a_\ell(3z+1).
\end{equation}
For every integer $H\geq1$ the actual sharp-arc error satisfies
\begin{equation}\label{A:ET-reduction}
 |R(X)-C_0X|\ll X^{2/3}+\frac XH+
       \sum_{1\leq\ell\leq H}\frac{|T_\ell(X)|}{\ell}
       +\log(2X).
\end{equation}
In particular, the uniform bound
\begin{equation}\label{A:needed-twist}
 |T_\ell(X)|\ll_\epsilon X^{2/3+\epsilon}
           +X^{1/2+\epsilon}\ell^{1/2}
 \quad(1\leq\ell\leq X^{1/3})
\end{equation}
would prove the requested exponent $2/3+\epsilon$.
We have not proved \eqref{A:needed-twist}.
\begin{proof}
Changing the generator by $\pm\varepsilon_0^k$ changes the logarithm
of the absolute embedding ratio by $2kL_0$, and so leaves
$\psi_\ell$ unchanged. The character is multiplicative. On the
positive-norm generator of $m_z$ it is exactly
$e^{2\pi i\ell\theta}$. For the unique prime ideal above $2$,
$|(1+\sqrt3)/(1-\sqrt3)|=\varepsilon_0$, so its character is
$(-1)^\ell$. Away from $2$, the integers $3z-1$ and $3z+1$ are
coprime. At $2$, where they can share powers, there is just one prime
ideal, and its local coefficient is geometric. Combining these facts
proves \eqref{A:twisted-shift}, including the ramified factor.
Conjugation pairs inverse character values, so $a_\ell(m)$ is real.

Apply the Erd\H{o}s--Tur\'an inequality in the form of
\cite[equation (8.2)]{CarneiroVaaler2010} to the finite list of the
$S(X)$ ideal angles. That equation bounds the translate of the sawtooth
sum; the characteristic function of an interval minus its length is
the difference of two translates of that sawtooth function, away from
its endpoints. Thus the arc discrepancy is at most a fixed constant
times $S(X)/H+\sum_{\ell\leq H}|T_\ell(X)|/\ell$, plus endpoint
multiplicities. The upper endpoint is empty by the calculation in
Statement A.4.1. At the lower endpoint, $t=0$ gives
$s^2-2(3z)^2=-2$. Reduction in $\mathbb Z[\sqrt2]$ by its norm-one
unit $3+2\sqrt2$ puts a positive generator in
$[\sqrt2,(3+2\sqrt2)\sqrt2)$, where the only solutions of norm $-2$
are the finite seeds $(s,3z)=(0,1),(4,3)$. Thus the total number with
$z\leq X$ is $O(\log(2X))$; imposing divisibility of the second
coordinate by $3$ can only decrease it.
Statement A.3.1 gives $S(X)=cX+O(X^{2/3})$ and $S(X)\ll X$.
Statement A.1.3 now proves \eqref{A:ET-reduction}.
Under \eqref{A:needed-twist}, take $H=\lfloor X^{1/3}\rfloor$;
the two sums cost $O_\epsilon(X^{2/3+\epsilon}\log X)$ and
$O_\epsilon(X^{1/2+\epsilon}H^{1/2})$, respectively. Relabeling
$\epsilon$ proves the asserted conditional implication.
\end{proof}

\paragraph{Uniformity audit and current obstruction.}
A bound for each fixed $\ell$ with an unspecified constant depending on
$\ell$ cannot be inserted in \eqref{A:ET-reduction} with $H$ growing
with $X$. More quantitatively, a bound
$T_\ell(X)\ll_\epsilon \ell^A X^{\beta+\epsilon}$, $A>0$, leads by this
argument to exponent $(A+\beta)/(A+1)$, after balancing $X/H$ against
$X^\beta H^A$. Thus even a fixed-mode exponent $2/3$ with polynomial
loss in the mode does not establish the sharp-arc exponent $2/3$.
Theorems 1.2 and 1.5 of Lekkas--Petridis integrate the entire source
geodesic. Their Remark 6 changes the target geodesic, but does not
supply these growing-mode estimates. Voskou's four orientation cases
\cite{Voskou2024} are discrete sign cases, not intervals in the foot
parameter. The fixed compact angular modes in
\cite[Theorem 1.1]{BourgainKontorovichSarnak2010} likewise do not
supply the required uniform hyperbolic-foot estimate.
The direct domain method of Statement A.4.1 proves $7/8+\epsilon$.
A joint radial and angular spectral smoothing with the uniform
estimate \eqref{A:needed-twist}, or a substitute averaged estimate
strong enough for \eqref{A:ET-reduction}, remains the concrete step
needed to reach $2/3+\epsilon$.
